% !TeX program = lualatex
% !TeX encoding = utf8
% !TeX spellcheck = uk_UA
% !TeX root =../PyscfBook.tex

%=========================================================
%\Opensolutionfile{answer}[\currfilebase/\currfilebase-Answers]
\chapter{Теорія відгуку та властивості молекул}\label{\currfilebase}
%=========================================================

%% ========================================================
\section{Теорія відгуку: загальна концепція}
%% ========================================================

Розглянемо систему, описану гамільтоніаном $\hat{H}_0$, яка перебуває в стаціонарному стані $\Psi_0$ з енергією $E_0$.
Якщо на неї діє зовнішнє збурення $\lambda \hat{V}$, то повний гамільтоніан набуває вигляду
\[
\hat{H}(\lambda) = \hat{H}_0 + \lambda \hat{V}.
\]
Параметр $\lambda$ визначає \emph{силу збурення}: це може бути електричне або магнітне поле, деформація геометрії, зміна спіну тощо.
Відповідно, хвильова функція та енергія системи теж залежать від $\lambda$:
\[
\Psi(\lambda) = \Psi^{(0)} + \lambda\,\Psi^{(1)} + \tfrac{1}{2}\lambda^2\Psi^{(2)} + \cdots,
\]
\[
E(\lambda) = E^{(0)} + \lambda E^{(1)} + \tfrac{1}{2}\lambda^2 E^{(2)} + \cdots.
\]

%% --------------------------------------------------------
\subsection{Розклад енергії за збуренням}
%% --------------------------------------------------------

Підставимо цей розклад у рівняння Шредінґера
\[
\hat{H}(\lambda)\Psi(\lambda) = E(\lambda)\Psi(\lambda),
\]
і прирівняємо члени однакового порядку за $\lambda$.
Для нульового, першого й другого порядків отримаємо:
\[
\begin{aligned}
\hat{H}_0 \Psi^{(0)} &= E^{(0)} \Psi^{(0)}, \\[3pt]
(\hat{H}_0 - E^{(0)})\Psi^{(1)} &= -(\hat{V} - E^{(1)})\Psi^{(0)}, \\[3pt]
(\hat{H}_0 - E^{(0)})\Psi^{(2)} &= -(\hat{V} - E^{(1)})\Psi^{(1)} + E^{(2)}\Psi^{(0)}.
\end{aligned}
\]

А поправки до енергії виражаються через оператор збурення $\hat{V}$ як:
\[
E^{(1)} = \langle \Psi^{(0)} | \hat{V} | \Psi^{(0)} \rangle,
\qquad
E^{(2)} = \langle \Psi^{(0)} | \hat{V} | \Psi^{(1)} \rangle.
\]

Це --- основа \textbf{теорії збурень}. Покажемо, як з неї випливає теорія відгуку.

Оскільки енергію системи можна виразити як сподіване значення гамільтоніана системи:

\[
E(\lambda) = \frac{\langle \Psi(\lambda) | \hat{H}(\lambda) | \Psi(\lambda) \rangle}
{\langle \Psi(\lambda) | \Psi(\lambda) \rangle}.
\]

То згідно теореми \textbf{Гельмана-Фейнмана}, похідна енергії по параметру збурення $\lambda$ має вигляд:
\[
\frac{dE}{d\lambda}
= \bracket<\Psi(\lambda) | \frac{d\hat{H}}{d\lambda} | \Psi(\lambda) >.
\]


Якщо далі підставити $\hat{H}(\lambda) = \hat{H}_0 + \lambda \hat{V}$, то отримаємо:
\[
\left.\frac{dE}{d\lambda}\right|_{\lambda=0}
= \langle \Psi^{(0)} | \hat{V} | \Psi^{(0)} \rangle
\;\equiv\; E^{(1)}.
\]
Отже, $E^{(1)}$ --- це не просто коефіцієнт при $\lambda$, а \emph{справжня похідна енергії за збуренням}.
Аналогічно, друга похідна
\[
\frac{d^2E}{d\lambda^2} = 2\langle \Psi^{(0)} | \hat{V} | \Psi^{(1)} \rangle.
\]


Покажемо далі, що \emph{усі властивості молекули можна виразити через похідні енергії за параметром, який описує збурення системи.}

%% --------------------------------------------------------
\subsection{Від енергії до властивостей}
%% --------------------------------------------------------

Нехай параметр~$F$ позначає будь-яку \emph{зовнішню дію на систему}:
це може бути електричне або магнітне поле, зміна координат ядер, деформація потенціалу чи спінове збурення.
Тоді повна енергія молекули є аналітичною функцією цього параметра:
\[
E(\mathbf{F}) = E_0
+ \left( \frac{\partial E}{\partial F_\alpha} \right) F_\alpha
+ \frac{1}{2} \left( \frac{\partial^2 E}{\partial F_\alpha \partial F_\beta} \right) F_\alpha F_\beta
+ \cdots
\]
де $F_\alpha$ --- компоненти зовнішньої дії  на систему.

Відповідні похідні визначають фізичні властивості:
\[
\begin{aligned}
\mu_\alpha &= - \frac{\partial E}{\partial E_\alpha}
\quad &&\text{(дипольний момент)},\\[3pt]
\alpha_{\alpha\beta} &= - \frac{\partial^2 E}{\partial E_\alpha \partial E_\beta}
\quad &&\text{(поляризованість)},\\[3pt]
\sigma_{\alpha\beta} &= - \frac{\partial^2 E}{\partial B_\alpha \partial M_\beta}
\quad &&\text{(магнітне екранування)},\\[3pt]
F_{A\alpha} &= - \frac{\partial E}{\partial R_{A\alpha}}
\quad &&\text{(сила на ядрі~$A$)},\\[3pt]
H_{AB}^{\alpha\beta} &= \frac{\partial^2 E}{\partial R_{A\alpha}\partial R_{B\beta}}
\quad &&\text{(гессіан, що визначає частоти коливань)}.
\end{aligned}
\]

Таким чином, як електричні, магнітні, так і механічні властивості молекули
описуються єдиним математичним підходом — через похідні енергії за параметрами, що збурюють систему.
Власне, \textbf{теорія відгуку} --- це узагальнення теорії збурень, в якій
замість формального розкладу хвильової функції розглядається її \emph{варіаційний відгук} на зовнішнє поле.


\begin{table}[h!]
\centering
\caption{Зв’язок порядку похідних енергії з молекулярними властивостями. Тут $q$ --- координати, $E$ --- зовнішнє електричне поле, $B$ --- зовнішнє магнітне поле, $I$ --- магнітне поле ядра.}\label{tab:prop}
\begin{tblr}{
%  cells = {font=\scriptsize},
  colspec={Q[c,m ]Q[c,m] Q[c,m] Q[c,m] Q[l,m]},
  cell{1}{1} = {c=4}{c},
  cell{1}{5} = {}{c,m},
  cell{2-Z}{1-4}={}{fg=blue, cmd=\bfseries},
  stretch = -1,
  hline{2} = {2pt},
  hline{3} = {1pt},
  hline{Z} = {1pt},
}
$n$-та похідна енергії &       &       &       & {Збурення} \\
$n_R$ & $n_E$ & $n_B$ & $n_I$ & Властивість \\
1 & 0 & 0 & 0 & Сила \\
2 & 0 & 0 & 0 & Частоти гармонічних коливань \\
3 & 0 & 0 & 0 & Ангармонічні поправки \\
0 & 1 & 0 & 0 & Електричний дипольний момент $\vec{\mu}_e$ \\
0 & 2 & 0 & 0 & Електрична поляризованість $\chi_e$ \\
0 & 3 & 0 & 0 & Електрична надполяризованість \\
0 & 0 & 1 & 0 & Магнітна сприйнятливість $\chi_m$ \\
0 & 0 & 2 & 1 & Константа надтонкого зв’язку (\texttt{EPR}) \\
0 & 0 & 0 & 2 & Стала спін-спінової взаємодії ядер \\
1 & 1 & 0 & 0 & Інтенсивності ІЧ-спектрів \\
2 & 1 & 0 & 0 & Інтенсивності обертонів в ІЧ-спектрах \\
1 & 2 & 0 & 0 & Інтенсивності Раманівських спектрів \\
2 & 2 & 0 & 0 & Інтенсивності обертонів у Раманівських спектрах \\
0 & 1 & 1 & 0 & Циркулярний дихроїзм \\
0 & 2 & 1 & 0 & Магнітний циркулярний дихроїзм \\
0 & 0 & 1 & 1 & Хімічний зсув (\texttt{NMR}) \\
\end{tblr}
\end{table}

Концепція теорії відгуку надзвичайно універсальна: практично всі молекулярні властивості можна розглядати як похідні енергії за зовнішніми параметрами --- електричним або магнітним полем, координатами ядер тощо.
Нижче наведено узагальнену схему, що показує, як похідні різних порядків визначають ті чи інші властивості (див. табл \ref{tab:prop}).

Як видно, усі ці властивості --- від сил та частот до магнітного екранування ---
описуються похідними енергії різних порядків за різними збуреннями.
Таким чином, \emph{усі властивості молекули є лише різними проявами одного й того ж принципу: відгуку хвильової функції на зовнішнє збурення.}

%% --------------------------------------------------------
\subsection{Формалізм варіацій орбіталей}
%% --------------------------------------------------------

У загальній теорії відгуку рівняння Шредінґера описує зміну повної хвильової функції~$\Psi(\lambda)$ під дією збурення.
Однак у багаточастинкових системах (молекулах) така функція залежить від великої кількості координат електронів,
і безпосередньо працювати з нею практично неможливо.

У методі Хартрі–Фока ця складна залежність замінюється набором одноелектронних функцій --- молекулярних орбіталей~$\phi_i$,
з яких будується детермінант Слейтера.
Саме ці орбіталі і є основними варіаційними змінними, що визначають енергію системи.

Тому, щоб реалізувати теорію відгуку на практиці, необхідно перейти від рівняння Шредінґера до орбітального формалізму,
у якому збурення описується через варіації орбіталей, а зміна енергії --- через відповідні зміни операторів Фока.

У методі Хартрі-Фока хвильова функція є детермінантом із молекулярних орбіталей~$\phi_i$, що задовольняють рівняння
\[
\hat{F}\,\phi_i = \varepsilon_i \phi_i.
\]
При прикладенні зовнішнього поля $\lambda \hat{V}$ орбіталі змінюються:
\[
\phi_i \;\to\; \phi_i + \lambda\,\phi_i^{(1)} + \mathcal{O}(\lambda^2),
\]
де $\phi_i^{(1)}$ --- \emph{варіації орбіталей}, або \textbf{збурений відгук} системи.
Підставивши це у рівняння Хартрі---Фока й розклавши за степенями $\lambda$, отримаємо лінійні рівняння для відгуку:
\[
(\hat{F}-\varepsilon_i)\phi_i^{(1)} = -\hat{F}^{(1)}\phi_i,
\]
які відомі як \textbf{рівняння CPHF} (Coupled-Perturbed Hartree--Fock).

Розв’язання цих рівнянь дозволяє визначити, як змінюються орбіталі під дією зовнішнього поля,
а отже --- обчислити похідні енергії, тобто молекулярні властивості.
У методах на основі функціонала густини (\texttt{DFT}) аналогічні рівняння мають назву \textbf{CPKS} (Coupled–Perturbed Kohn--Sham).

На практиці цей формалізм реалізовано в пакеті~\texttt{PySCF},
де користувач може вибрати, чи враховувати релаксацію орбіталей у відгуку:
\begin{itemize}
    \item \inlinecode{cphf=False} --- обчислення без урахування відгуку орбіталей (\emph{unrelaxed} властивості);
    система не адаптується до збурення.
    \item \inlinecode{cphf=True} --- повне самозгоджене врахування відгуку (\emph{relaxed} властивості);
    орбіталі релаксують під дією поля, що забезпечує фізично коректні похідні енергії.
\end{itemize}

\input{\currfiledir ElectricProperties.tex}


